\begin{papercard}{mahatody2010}
\cf{Stages} \stagetag{W}. A state-of-the-art review of the cognitive walkthrough: the three original versions (CW1--CW3) and eleven variants, with four syntheses (concepts, comparative studies, non-comparative studies, choice of variant). \emph{Note:} the appended copy is the author manuscript, so its pages differ from the journal pagination (IJHCI 26(8), 741--785).

\cf{Definition} CW is a usability inspection method that links the interface walkthrough to a cognitive model. The evaluator performs typical tasks, and actions and responses are judged against the user's goals and knowledge \pc{mahatody2010}{2}{Cognitive Walkthrough (CW) is a usability inspection method that links the interface walkthrough to a cognitive model}. Common principle of all versions: simulate the user's cognitive behaviour by answering questions tied to a cognitive model \pc{mahatody2010}{2}{the method simulates the cognitive behavior of the user by responding to questions}.

\cf{CW1 (Lewis et al.\ 1990)} For walk-up-and-use interfaces, based on CE+ exploratory learning \pc{mahatody2010}{2}{was designed to evaluate walk- up-and-use interfaces}. Two phases: preparation (tasks broken into action sequences; evaluators may be designers, users or experts) and evaluation (questions per action). It identified only 50\,\% of observed errors \pc{mahatody2010}{3}{noted that the method was promising but not yet satisfactory because only 50}.

\cf{CW2 (Polson et al.\ 1992)} Adds Norman's seven-stage action theory and Kintsch's construction--integration model. The user builds a goal structure similar to GOMS \pc{mahatody2010}{3}{the user must build a goal structure, which is similar to the hierarchical structure}. Three forms: goal match, action selection, goal update after feedback. More cumbersome than CW1, which led to an automated HyperCard version (AutoCW2) \pc{mahatody2010}{3}{According to its authors, this version is more complex and cumbersome to use than CW1}, \pc{mahatody2010}{4}{proposed an automated version (AutoCW2)}.

\cf{CW3 (Wharton et al.\ 1994)} A credible scenario per action, justified by the user's background and the interface feedback, with four questions: right effect, notice the action, associate action with effect, see progress \pc{mahatody2010}{4}{In CW3, the evaluator is invited to imagine a specific and credible scenario}. Reviewers agree it is easy to learn but tedious \pc{mahatody2010}{4}{Most of the reviewing authors agreed that CW is tedious}.

\cf{Eleven variants} \emph{Heuristic Walkthrough} (Sears 1997): phase 1 uses prioritised tasks and the CW3 questions, phase 2 free exploration with heuristics. It found more problems than CW3 and fewer false positives than HE \pc{mahatody2010}{4}{Sears showed that HW can find more usability problems than CW3}. \emph{Norman CW} (Rizzo et al.\ 1997): six questions on cognitive distance for distributed, multicultural teams \pc{mahatody2010}{5}{none of the initial 3 versions of Cognitive Walkthough is particularly appropriate}. \emph{Streamlined CW} (Spencer 2000): five phases, two questions and ground rules \pc{mahatody2010}{6}{SCW is a variant of CW3 with 5 phases}. \emph{CW for the Web} (Blackmon et al.\ 2002): the CoLiDeS model and Latent Semantic Analysis replace subjective judgement. Goals are written in 100--200 words. A heading with LSA below 0.8 counts as unfamiliar, and a pair above 0.6 as confusable \pc{mahatody2010}{7}{a title with an LSA value of less than 0.8 is assumed to be unfamiliar}. \emph{Groupware Walkthrough} (Pinelle \& Gutwin 2002): task model with all paths, and effectiveness, efficiency and satisfaction questions after each task \pc{mahatody2010}{8}{Every possible path is explored during the walkthrough}. \emph{Activity Walkthrough} (Bertelsen 2004): six phases with contextualisation from activity theory \pc{mahatody2010}{9}{In Activity Walkthrough, the evaluation is conducted in six phases}. \emph{Interaction Walkthrough} (Ryu \& Monk 2004): goal--action, action--effect and effect--goal problems, including hidden or misleading modes \pc{mahatody2010}{10}{IW evaluation looks at the three paths of the interaction cycle}. \emph{CW with Users} (Granollers \& Lorés 2006): expert CW, then think-aloud users, then expert review of users' doubts \pc{mahatody2010}{10}{Cognitive Walkthrough with Users (Granollers}. \emph{Extended CW} (Kato \& Hori 2005): nine questions separating object from action and perceivability from understandability, aimed at accessibility \pc{mahatody2010}{11}{they included questions Q2-a and Q2-c for assessing perceptibility}. \emph{Distributed CW} (Eden 2007): four questions on how represented information supports the task \pc{mahatody2010}{12}{Q1: Will the way that information is represented provide all the knowledge required}. \emph{Enhanced CW} (Bligård \& Osvalder 2007): nine questions rated 1--5 with justification, where anything below 5 is a potential problem, plus five problem categories \pc{mahatody2010}{12}{The answer to each question is a number ranging between 1 and 5}, \pc{mahatody2010}{13}{Any answer other than 5 suggests a potential usability problem}.

\cf{Common core} Across all variants the questions ask whether users (i) formulate goals, (ii) perform the appropriate actions, (iii) interpret system responses correctly and (iv) accomplish their goals \pc{mahatody2010}{15}{in general the questions are designed to evaluate, for each interaction}.

\cf{Evidence from comparative studies} CW1: 20 problems, 18 found by at least two of four evaluators, and 50\,\% of empirical problems \pc{mahatody2010}{16}{The authors highlighted 20 problems, 18 of which were found by at least two}. CW2: 30\,\% of problems but 70\,\% of serious ones on a text editor \pc{mahatody2010}{16}{CW2 detected 30\% of the usability problems, including 70\% of the problems considered to be serious}. Jeffries et al.\ (1991, on HP-VUE): HE found 105 problems, usability testing 31, guidelines 35 and CW2 35. CW2 can be used by non-specialist developers \pc{mahatody2010}{18}{The heuristic evaluation method allowed 105 usability problems to be identified}. John \& Marks (1997): CW3 found 42 problems with only 5\,\% false alarms (HE 17\,\%), and 73\,\% of its implemented-aspect problems led to code changes \pc{mahatody2010}{18}{However, 2 (5\%) of the 42 problems found by CW3 were false alarms}. Huart et al.\ (2004): effectiveness depends heavily on evaluator expertise \pc{mahatody2010}{20}{the effectiveness of the method depends greatly on the expertise of the evaluator}.

\cf{Conclusions} No study helps choose among variants \pc{mahatody2010}{21}{We did not find any studies intended to help the evaluators}. Validation is thin: few tests, often subjective, and no common criteria \pc{mahatody2010}{21}{No comprehensive study has ever been made to compare the different versions and extensions}. Methodological weaknesses remain in preparation (evaluator training, task analysis, context) and in using the results. The authors recommend audio-video recording \pc{mahatody2010}{22}{Many methodological weaknesses remain in terms of preparing for the assessment}. Many domain adaptations exist, including medical \pc{mahatody2010}{22}{have adapted certain versions and variants of CW for evaluating HCI in the medical domain}.

\cf{Use for this thesis} A map for \emph{choosing and naming} the CW variant. A privacy-information walkthrough on apps resembles CWW (finding information through labels and headings) and Enhanced CW (graded answers with justifications and problem categories). Graded 1--5 answers with justification could make the walkthrough codeable for two-coder agreement. Cite this review when stating that CW variants are under-validated and that any adaptation must be described explicitly, e.g.\ ``CW3 questions adapted to privacy-information tasks''. The evidence that results depend on expertise supports reporting the inspector's background.
\end{papercard}
